% Generated draft: requires mathematical and source review.
\documentclass[11pt]{article}
\usepackage[margin=25mm]{geometry}
\usepackage{fontspec}
\setmainfont{Latin Modern Roman}
\usepackage{amsmath,amssymb,mathtools,unicode-math}
\setmathfont{Latin Modern Math}
\usepackage{enumitem,needspace,xurl,hyperref}
\setlength{\parindent}{0pt}
\setlength{\parskip}{5pt}
\setlength{\emergencystretch}{3em}
\newcommand{\N}{\mathbb{N}}
\newcommand{\Z}{\mathbb{Z}}
\newcommand{\Q}{\mathbb{Q}}
\newcommand{\R}{\mathbb{R}}
\newcommand{\C}{\mathbb{C}}
\newcommand{\F}{\mathbb{F}}
\newcommand{\A}{\mathbb{A}}
\newcommand{\PP}{\mathbb{P}}
\newcommand{\E}{\mathbb{E}}
\newcommand{\eps}{\varepsilon}
\newcommand{\dd}{\,\mathrm{d}}
\begin{document}
\section*{UnsolvedMath ID: 11000314}
\textbf{Source problem code:} catalogue entry
\subsection{Definitions and mathematical statement}
Problem 4.8 asks to define unstable homology and/or cohomology classes of \(\mathrm{GL}(n,\Z)\). The term unstable is not defined: it could refer to classes outside a stability range, classes killed by stabilization, or another rank-dependent construction. These meanings are not automatically equivalent. The standard group inclusion is \(A\mapsto\operatorname{diag}(A,1)\), inducing homology forward and cohomology backward, but the source gives no chosen criterion for instability.

Its particular question concerns source-defined classes
\[
 \mu_i\in H_{4i}(\operatorname{Out}(F_{2i+2});\Q).
\]
The abelianization action induces a natural homomorphism
\[
 a_i:\operatorname{Out}(F_{2i+2})
       \longrightarrow\mathrm{GL}(2i+2,\Z).
\]
The requested image, with matching indices, is
\[
 (a_i)_*(\mu_i)\in H_{4i}(\mathrm{GL}(2i+2,\Z);\Q).
\]
This is a homological image, not a cohomological pullback. The source prints the projection using \(k\) instead of \(i\), namely rank \(2k+2\); it must take \(k=i\) for the displayed class to be in that map's domain unless additional missing comparison maps are supplied. That notational mismatch is explicit here.

The construction, index range, normalization and known properties of \(\mu_i\) are absent. The source asks what can be said about the images, not solely whether one selected image is nonzero. It says preceding results and computations support a following possibly folklore conjecture, but neither those computations nor that conjecture is included, and no guessed nonvanishing or vanishing statement is inserted.
\subsection{Short English statement}
Construct unstable classes for integer linear groups and analyze the specified rational homology images from free-group outer automorphism groups. Preserve degree \(4i\) and rank \(2i+2\), while flagging the source's mismatched \(k\)-index and missing definition of the classes and instability.
\subsection{Sources}
\begin{itemize}
\item Local UnsolvedMath registry, permanent ID \(11000314\), exact source code \texttt{catalogue entry}, Problem 4.8, PDF page 366. No individual proposer or proposed year is supplied.
\item Farb, edited volume \emph{Problems on Mapping Class Groups and Related Topics}: \url{https://www.math.uchicago.edu/~farb/papers/mcgbook.pdf}. The volume has 329 explicitly typeset numbered items across chapters, repeated local numbering distinguished by PDF page.
\item No complete reference for \(\mu_i\), the preceding computations, or the following possible-folklore conjecture is supplied in the record; none is reconstructed.
\item Extraction pdf-font-structure, accessed 2026-07-29. Source evidence NEEDS\_REVIEW presents an open problem and requires a current primary-literature audit. Supplied OPEN-TRIAGE summary describes an unstable-linear-group-(co)homology construction program with partial progress, without precise cases or citations.
\item Dataset attribution: Ulam AI / UnsolvedMath Contributors, CC BY 4.0. Source rights NEEDS\_REVIEW: publicly accessible author/publisher copy, redistribution terms requiring review.
\end{itemize}
\end{document}
